\section*{Overview}

Many contest problems are not fundamentally about geometry or fancy data structures. They are about noticing that all
important changes happen at a finite set of coordinates. Once those coordinates are sorted, the problem becomes a sweep
over \emph{events}.

\section*{When to Use It}

This pattern is a good fit when:

\begin{itemize}
  \item intervals start and end,
  \item queries ask ``what is active at coordinate \(x\)?'',
  \item the answer only changes when the sweep line crosses an endpoint, a query point, or another special coordinate.
\end{itemize}

It is one of the cleanest ways to reduce \(O(nq)\) scans into \(O((n+q)\log(n+q))\) offline processing.

\section*{Core Idea}

Encode each structural change as an event:

\begin{itemize}
  \item add an interval,
  \item remove an interval,
  \item answer a query.
\end{itemize}

Sort all events by coordinate, decide a tie-breaking order, and maintain the active state while sweeping from left to
right.

\section*{Key Insight}

The real invariant is not ``where the sweep line is.'' It is ``which objects intersect the current position.'' If that
active set can be updated locally when crossing one event, then the entire problem becomes a sequence of small changes
instead of repeated full recomputation.

\section*{Worked Problem}

\paragraph{Problem.}
You are given closed intervals \([l_i, r_i]\) and query points \(x_j\). For each query, count how many intervals
contain it.

\paragraph{Why events fit.}
An interval contributes \(+1\) starting at \(l_i\) and stops contributing after \(r_i\). A query only needs the number
of intervals that are active exactly at its coordinate.

\paragraph{Algorithm outline.}

\begin{itemize}
  \item Create events \((l_i, \texttt{add})\), \((r_i, \texttt{remove})\), and \((x_j, \texttt{query})\).
  \item Sort by coordinate.
  \item For equal coordinates, process \texttt{add} before \texttt{query} before \texttt{remove} for closed intervals.
  \item Maintain one integer \texttt{active}.
\end{itemize}

\section*{Problem Pattern}

This exact structure also appears in:

\begin{itemize}
  \item maximum overlap of intervals,
  \item offline stabbing queries,
  \item rectangle union and line sweeps, after the state becomes richer than one counter,
  \item counting how many resources are currently occupied.
\end{itemize}

\section*{Correctness Intuition}

Between two consecutive event coordinates, the active set does not change. So it is enough to update the state only at
events. The tie-breaking rule is the only place where the interval semantics matter: closed, open, and half-open
intervals need different ordering.

\section*{Complexity Analysis}

If there are \(E\) total events, sorting costs \(O(E \log E)\) and the sweep is linear after sorting.

\section*{Implementation}

The sample code solves the interval stabbing problem above. The same event representation extends naturally when the
active state becomes a Fenwick tree, multiset, or segment tree instead of one integer.

\section*{Common Pitfalls}

\begin{itemize}
  \item Getting tie-breaking wrong at equal coordinates.
  \item Mixing closed-interval logic with half-open implementation.
  \item Forgetting that offline event sorting loses the original query order unless IDs are stored.
  \item Using an event sweep when the hard part is really a graph interaction, not a one-dimensional ordering.
\end{itemize}

\section*{Variants / Extensions}

\begin{itemize}
  \item Coordinate compression before the sweep.
  \item Two-dimensional sweep line with an auxiliary data structure.
  \item Difference-array interpretation when events only change one counter.
\end{itemize}

\section*{Practice Problems}

\begin{itemize}
  \item Count interval coverage at query points.
  \item Find the maximum number of simultaneously active segments.
  \item Offline process add/remove/query actions sorted by one key.
\end{itemize}

\section*{References}

\begin{itemize}
  \item \href{https://usaco.guide/plat/sweep-line}{USACO Guide: Sweep Line}.
  \item \href{https://codeforces.com/catalog}{Codeforces Catalog}.
  \item \href{https://en.oi-wiki.org/geometry/scanning/}{OI Wiki: Scanning Line}.
\end{itemize}
